% !TEX root = analysis_class-note.tex
\chapter{Sep 30}
\clearpage
\section{Sep 30}
\dfn[archimedean-property]{The Archimedean Property}{
    An ordered field $\mcF$ has the Archimedean property if, for every $x \in \mcF$, there exists an $n \in \bbN$ such that $x < n$.
}

\cor[lub-implies-archimedean]{}{
    If $\mcF$ is an ordered field with the least upper bound property, then $\mcF$ is Archimedean.
}
\begin{proof}
    Suppose, for the sake of contradiction, that $\bbN$ is bounded above in $\mcF$. Since $\mcF$ has the least upper bound property, we can let $\alpha = \sup \bbN$. Then $\alpha \ge n$ for all $n \in \bbN$. Since $n+1 \in \bbN$, we also have $\alpha \ge n+1$, so $\alpha-1 \ge n$ for all $n \in \bbN$. Thus $\alpha-1$ is an upper bound, but $\alpha-1 < \alpha$, contradicting the fact that $\alpha$ is the least upper bound.
\end{proof}

\lem[density-of-rationals]{}{
    For all $x,y \in \bbR$ with $x < y$, there exists $q \in \bbQ$ such that $x < q < y$.
}

\begin{proof}
    Since $y-x > 0$, we have $\frac{1}{y-x} > 0$. By the Archimedean property, there exists an $n \in \bbN$ such that $\frac{1}{y-x} < n$, or equivalently, $\frac{1}{n} < y-x$.

    Now choose $m \in \bbZ$ such that $m-1 \le nx < m$ (that is, $m = \min\{k \in \bbZ : k > nx\}$). Then
    \[
        x < \frac{m}{n} \le x + \frac{1}{n} < y.
    \]
    Taking $q = \frac{m}{n} \in \bbQ$ gives the result.
\end{proof}

\thm[small-reciprocals]{}{
    For any $\veps > 0$, there exists $n \in \bbN$ such that $0 < \frac{1}{n} < \veps$.
} 
\begin{proof}
    By the Archimedean property, there exists $n \in \bbN$ such that $n > \frac{1}{\veps}$. Hence $0 < \frac{1}{n} < \veps$.
\end{proof}

\dfn[function-rule]{Functions}{
    A real-valued function of a real variable is a ``rule'' that assigns a real number to each real number in some set.
}

\begin{enumerate}[label=(\alph*)]
    \item A function is denoted by $f$, and $f(x)$ is the value assigned to $x$ by $f$.
    \item The domain of $f$ is the set of all real numbers to which the rule applies.
\end{enumerate}

\dfn[function-graph]{}{
    A real-valued function of a real variable can be represented by a subset $\mcF \subseteq \bbR \times \bbR$ with the following properties:
    \begin{enumerate}[label=(\alph*)]
        \item If $(a,b)$ and $(a,c)$ are in $\mcF$, then $b=c$.
        \item The domain of the function is the set
        \[
            D = \{a \in \bbR : (a,b) \in \mcF \text{ for some } b \in \bbR\}.
        \]
        By (a), this $b$ is unique for each $a \in D$, and we write $b=f(a)$.
    \end{enumerate}
}


\section{Oct 2}

\dfn[key]{functions}{
    A function $f: S \to J$ is a subset \[
        F \subseteq S \times J = \{ (s,t ) : s \in S, t \in J \}
    \]
    s.t. if $(a, b ) \in F$ and $(a,c) \in F$ then $b = c$ we write $f(a) = b$      
}
Terminology
\begin{enumerate}[label=(\alph*)]
    \item The Domain of $f $ is denoted $Dom(f) $ is defined by \[
        Dom(F) = \{ a \in S : (a,b) \in F \text{ for some } b \in \bbR \}
    \]
     
    \item The Codomain of $F$ is $J$
    \item The Range if $f$ is \[
        Ran(F) = \{ b \in J : (a, b )  \in F \text{ for some } a \in S \}
    \]
       
\end{enumerate}

Some graphs...

\rmk[key]{}{
    \begin{enumerate}[label=(\alph*)]
        \item $F$ is injective if, for each $b \in Ran(F)  $  there is a unique $a \in Dom(F )$ s.t $(a, b ) \in F$  
        \item $F$ is surjective if $Ran(F) = J$
        \item $F$ is bijective if $F$ is injective and surjective.    
    \end{enumerate}
}

\dfn[key]{}{
    If $f: S \to J$ is a function, then we can always define \[
        f ^{-1 } = \{ (b , a ) \in J \times S : (a,b) \in F \subseteq S \times J \} 
    \]
     
}

Some graph...

\thm[key]{}{
    $f ^{-1}$ is a function $\iff $f is injective.
}

\begin{proof}
    Need to show if $(b_1, a_1)$ and $(b_1. a_1)$ in $f ^{-1}$ and $b_1 = b_2 \Rightarrow a_1 = a_2$
    $\iff f(a_1) = f(a_2) \text{ then } a_1 = a_2 \iff f \text{ is injective }  $
\end{proof}
Even if $f: S \to J$ is not injective, we can always define the \textbf{inverse image} of a set $B \subseteq J$ as follow, 
\[
    f ^{-1} (B) = \{ a \in S : (a, b) \in F \text{ for some } b \in B \}
\]
 

\[
    F = \{ (a, f(a) ) :a \in Dom(f)\}
\]

\section{compositions of functions}
\dfn[key]{compositions}{
    If $F: J_1 \times J_2$  and $G: J_0 \times J_1$ and function can define $F \circ G: J_0 \times J_2$ by $$f \circ g(a) = f(g(a))$$
    Also,
    \[
        Dom(f \circ g) = \{ a \in Dom(g) : g(a) \in Dom(f) \}
    \]
    
}

Functions whose domain is $\bbN$
\dfn[key]{}{
    An infinite sequence $\{ X_n \}_{n \in \bbN} \subseteq J$ is a function $f: \bbN \to T$ with $Dom(f) = \bbN$, i.e. $x_n = f(n)$  
}
 
Think $X_n = \frac{1 }{n } $ is a sequence in $\bbR$.
\begin{enumerate}[label=(\alph*)]
    \item $X_n = sin(n)$ is a sequence in $\bbR$ 
    \item $X_n = n ^{th }$ digit of $\pi$  
\end{enumerate}
We often write $x_1, x_2, x_3, \dots$
A key Role in Mathematics is played by the notion of a convergent sequence.

\dfn[key]{}{
    Let $\{ X_n  \}_{n\in \bbN} \subseteq \bbR$ be a sequence of real numbers. We say that $\{ X_n \}$ converges to $L$, and write \[
        \lim_{n \to \infty} X_n = L
    \] 
    For every $\veps > 0$, there is a $N \in \bbN$, s.t. for all $n \ge \bbN$ we have, \[
        |X_n - L | < \veps
    \]
      Say $\{ X_n  \} $ converges if it converges to some limit.
}

%  \ex[key]{Archimedean}{
%     $\lim_{n \to \infty} \frac{1 }{n}$.
%     \[

%     \]
     
%  }
 
